% TOP_PROBLEM: {"id": 10000023, "title": "Self-avoiding loops on nonamenable transitive graphs", "category_id": 19, "category": {"id": 19, "name": "probability", "display_name": "Probability", "description": "Problems involving probability theory, stochastic processes, and random structures.", "slug": "probability", "order_index": 19, "created_at": "2026-07-31T15:26:25.671Z"}, "status": "open", "problem_number": "AMR-099-0023", "set_id": 13, "statusQualification": "Uses limsup for loop growth, avoiding the source’s undefined ordinary limit on bipartite graphs and trees."}
% FIRST_POSTED: 2026-10-01
% ENTRY_KIND: statement_only
% UnsolvedMath ID 10000023; source catalogue code AMR-099-0023
% !TeX program = lualatex
% Definition and sources only; this file is not an LLM solution attempt.
\documentclass[11pt,a4paper]{article}
\usepackage[margin=25mm]{geometry}
\usepackage{fontspec}
\setmainfont{lmroman10-regular.otf}[BoldFont=lmroman10-bold.otf,
 ItalicFont=lmroman10-italic.otf,BoldItalicFont=lmroman10-bolditalic.otf]
\usepackage{amsmath,amssymb,mathtools}
\usepackage{unicode-math}
\setmathfont{latinmodern-math.otf}
\AtBeginDocument{\let\setminus\smallsetminus}
\usepackage{xurl,hyperref}
\hypersetup{colorlinks=true,urlcolor=blue}
\setcounter{secnumdepth}{0}
\setlength{\parindent}{0pt}
\setlength{\parskip}{5pt}
\setlength{\emergencystretch}{3em}
\begin{document}
\section*{Self-avoiding loops on nonamenable transitive graphs}\label{10000023}
{\small UnsolvedMath ID 10000023 · AMR-099-0023 · Probability · AMR Open Problem Lists}
\subsection{Definitions and mathematical statement}
Let $G$ be an infinite connected locally finite vertex-transitive simple graph. For finite nonempty $A\subset V(G)$, let $\partial A$ be the vertices outside $A$ adjacent to it, and suppose $h(G)=\inf_A|\partial A|/|A|>0$.

Fix a root $o$. Let $c_n$ count sequences $(v_0,\ldots,v_n)$ with $v_0=o$, adjacent consecutive vertices, and all vertices distinct. These are rooted self-avoiding walks of length $n$. Their connective constant is
\[
\mu(G)=\lim_{n\to\infty}c_n^{1/n},
\]
whose existence follows from submultiplicativity and transitivity. Let $\ell_n$ count rooted oriented simple cycles $(v_0,\ldots,v_n)$ with $v_0=v_n=o$, $n\ge3$, and $v_0,\ldots,v_{n-1}$ distinct. Define the loop growth rate by
\[
\mu_{\mathrm{loops}}(G)=\limsup_{n\to\infty}\ell_n^{1/n},
\]
with $0^{1/n}=0$. Prove that $\mu_{\mathrm{loops}}(G)<\mu(G)$.

The limsup is necessary to make the definition valid without restrictions on cycle lengths: bipartite graphs have no odd cycles, and trees have no cycles at all. Reversing the orientation or counting each cycle without a distinguished orientation changes counts by at most a constant and therefore leaves the exponential rate unchanged.

The assertion is a strict comparison of exponential growth, not simply the observation that a walk has more possible endpoints than a loop. The desired gap may depend on $G$; no uniform positive gap over all nonamenable transitive graphs is specified. Both rates use the same unweighted graph and the same edge-length convention.
\subsection{Short English statement}
On every nonamenable transitive graph, do self-avoiding loops grow exponentially more slowly than self-avoiding walks?
\subsection{Sources}
\begin{itemize}
\item[S1] ULAM.AI and the UnsolvedMath Contributors, supplied problems.json, record 10000023 (AMR-099-0023), AMR Open Problem Lists. Catalogue identity and source transcription; CC BY 4.0.
\url{https://huggingface.co/datasets/ulamai/UnsolvedMath}.
\item[S2] Benjamini - Coarse Geometry and Randomness (Saint-Flour notes), Conjecture 4.32, PDF page 38. Original source indexed by AMR-099-0023.
\url{https://www.wisdom.weizmann.ac.il/~itai/stflouraug24.pdf#page=38}.
\item[S3] I. Benjamini, Coarse Geometry and Randomness, primary Saint-Flour notes; the numbered question and its surrounding definitions.
\url{https://arquivo.pt/noFrame/replay/20201231041548id_/http://www.wisdom.weizmann.ac.il/~itai/stflouraug24.pdf}.
\end{itemize}
\end{document}
